% !TeX root = ../main.tex
% This original sample is used when local-materials/rubric.pdf is absent.
\clearpage
\begingroup
\thispagestyle{empty}
\definecolor{rubricCrimson}{HTML}{A00018}
\AddToShipoutPictureBG*{%
  \AtPageCenter{\makebox(0,0){%
    \rotatebox{54}{\resizebox{1.33\paperwidth}{!}{%
      \sffamily\bfseries\color{rubricCrimson!65!white}EXAMPLE}}}}}

\setlength{\parindent}{0pt}
\setlength{\parskip}{0pt}
\vspace*{8mm}
{\sffamily\bfseries\small\color{rubricCrimson}REPLACE BEFORE SUBMISSION\par}
\vspace{4mm}
{\fontsize{26}{31}\selectfont\bfseries Sample marking rubric\par}
\vspace{4mm}
These criteria and point values illustrate how a marking sheet can be laid out.
Use the rubric supplied for your assignment when preparing the final report.\par

\vspace{13mm}
\noindent\begin{tblr}{
  width=\linewidth,
  colspec={X[1.15,l] X[1.6,l] Q[c,wd=1.55cm] Q[c,wd=1.75cm]},
  column{1}={leftsep=0pt}, column{4}={rightsep=0pt},
  rows={m,rowsep=6mm}, row{1,Z}={font=\bfseries,rowsep=3mm},
  hline{1,Z}={0.8pt}, hline{2,3,4,5}={0.35pt},
}
  Criterion & Sample evidence & Possible & Awarded \\
  Criterion A & Example outline & 4 & \\
  Criterion B & Example explanation & 4 & \\
  Criterion C & Example illustration & 2 & \\
  Total & & 10 & \\
\end{tblr}\par

\vspace{16mm}
{\bfseries Marker comments}\par
\vspace{12mm}\hrule height 0.35pt
\vspace{12mm}\hrule height 0.35pt
\vspace{12mm}\hrule height 0.35pt

\vfill
{\small To replace this page, save the course PDF as
\texttt{local-materials/rubric.pdf} and recompile.\par}
\clearpage
\endgroup
